% Options for packages loaded elsewhere
\PassOptionsToPackage{unicode}{hyperref}
\PassOptionsToPackage{hyphens}{url}
\documentclass[
  11pt,
]{article}
\usepackage{xcolor}
\usepackage[margin=1in]{geometry}
\usepackage{amsmath,amssymb}
\setcounter{secnumdepth}{-\maxdimen} % remove section numbering
\usepackage{iftex}
\ifPDFTeX
  \usepackage[T1]{fontenc}
  \usepackage[utf8]{inputenc}
  \usepackage{textcomp} % provide euro and other symbols
\else % if luatex or xetex
  \usepackage{unicode-math} % this also loads fontspec
  \defaultfontfeatures{Scale=MatchLowercase}
  \defaultfontfeatures[\rmfamily]{Ligatures=TeX,Scale=1}
\fi
\usepackage{lmodern}
\ifPDFTeX\else
  % xetex/luatex font selection
\fi
% Use upquote if available, for straight quotes in verbatim environments
\IfFileExists{upquote.sty}{\usepackage{upquote}}{}
\IfFileExists{microtype.sty}{% use microtype if available
  \usepackage[]{microtype}
  \UseMicrotypeSet[protrusion]{basicmath} % disable protrusion for tt fonts
}{}
\makeatletter
\@ifundefined{KOMAClassName}{% if non-KOMA class
  \IfFileExists{parskip.sty}{%
    \usepackage{parskip}
  }{% else
    \setlength{\parindent}{0pt}
    \setlength{\parskip}{6pt plus 2pt minus 1pt}}
}{% if KOMA class
  \KOMAoptions{parskip=half}}
\makeatother
\setlength{\emergencystretch}{3em} % prevent overfull lines
\providecommand{\tightlist}{%
  \setlength{\itemsep}{0pt}\setlength{\parskip}{0pt}}
\usepackage{xurl}
\usepackage{microtype}
\usepackage{mathrsfs}
\emergencystretch=2em
\AtBeginDocument{\hypersetup{pdftitle={A nonzero quantum elimination correction on two adjacent Wilson plaquettes},pdfauthor={}}}
\usepackage{bookmark}
\IfFileExists{xurl.sty}{\usepackage{xurl}}{} % add URL line breaks if available
\urlstyle{same}
\hypersetup{
  hidelinks,
  pdfcreator={LaTeX via pandoc}}

\author{}
\date{}

\begin{document}

\section{A nonzero quantum elimination correction on two adjacent Wilson
plaquettes}\label{a-nonzero-quantum-elimination-correction-on-two-adjacent-wilson-plaquettes}

8 September 2026. This is an explicit finite interacting calculation,
not a four-dimensional continuum mass-gap result. It evaluates a
vacuum-dependent elimination correction on a genuine gauge-invariant
retained observable.

\subsection{1. Original coefficients, graph and
operator}\label{original-coefficients-graph-and-operator}

Take two adjacent square plaquettes, P and Q, with a single common edge
e, six other distinct edges, and open boundary. This finite cell complex
has seven physical links; no link outside this specified complex is
presumed to have been integrated without its interaction. Choose
orientations so their ordered holonomies, after cyclic changes of their
base vertices, are \[
 P=U_e A,\qquad Q=U_e^{-1}B,
\] where A and B are the ordered products of the respective three outer
links. Set \(W_P=\operatorname{tr}P\), \(W_Q=\operatorname{tr}Q\). These
are real for SU(2). Each Haar measure has total mass one; the total
configuration space is \(\mathcal Q=SU(2)^7\), with product measure
\(d\lambda\). Use the anti-Hermitian basis \(T_a=i\sigma_a/2\),
\(-2\operatorname{tr}T_aT_b=\delta_{ab}\), and the actual Hamiltonian \[
 H= -\kappa\sum_{f=1}^7\Delta_f^T
       +b(4-W_P-W_Q),\qquad \kappa>0,\quad b>0.
\] Every plaquette of this complex is retained. Its relation to the
original action convention \(c=-\operatorname{tr}/2\) is \[
 \Delta^c=4\Delta^T,\qquad
 \kappa=\frac{2g_{\rm YM}^2}{a},\quad
 b=\frac1{2g_{\rm YM}^2a},\quad
 \xi:=\frac b\kappa=\frac1{4g_{\rm YM}^4}.
\] Thus the small parameter below is a specified strong-coupling
parameter, not a removal of the interaction at the positive values under
study.

The compact elliptic operator H has a unique strictly positive smooth
ground state. One proof is the strictly positive heat kernel for the
product Laplacian and its Feynman--Kac formula with the bounded real
potential. Compactness gives a lowest eigenfunction; replacing it by its
absolute value does not increase its form, and positivity improvement
makes the lowest eigenspace one-dimensional. Gauge transformations
commute with H and preserve positivity, so this ground state is gauge
invariant. Write its unit-L2 representative as \(\psi_\xi\), its
eigenvalue as \(E_\xi\), and
\(d\mu_\xi=\rho_\xi d\lambda=\psi_\xi^2d\lambda\). The unit-L2
convention fixes the state only: all scalar factors below are retained
explicitly.

\subsection{2. Controlled first derivative of the actual
vacuum}\label{controlled-first-derivative-of-the-actual-vacuum}

Let \(H_0=-\sum_f\Delta_f^T\). The exact operator identity is \[
 H=\kappa[H_0-\xi(W_P+W_Q)]+4bI.
\] The scalar \(4bI\) is not discarded from H; it changes its ground
energy by exactly 4b and leaves its ground vector and excitation
differences unchanged. The constant function is the simple zero
eigenvector of H0, and the next eigenvalue on the full link Hilbert
space is 3/4, from the spin-1/2 Casimir on one link. On the complex
circle \(|z|=3/8\), the resolvent of H0 has norm at most 8/3. Since
\(|W_P+W_Q|\leq4\), its resolvent Neumann series converges uniformly
when \(|\xi|<3/32\). The Riesz integral therefore gives an analytic
rank-one eigenprojection near zero. For real sufficiently small xi this
is the ground projection; the remainder of the spectrum stays separated
by bounded perturbation.

Choose its analytic eigenvector \(\phi_\xi\) with
\(\int\phi_\xi d\lambda=1\). This fixes a representative, not the
physical state norm. Put \[
 Z_\xi=\int\phi_\xi^2d\lambda,\qquad
 \psi_\xi=\phi_\xi/Z_\xi^{1/2},\qquad
 \rho_\xi=\phi_\xi^2/Z_\xi.
\] Haar integration of a link appearing once gives
\(\int W_P=\int W_Q=0\). Each plaquette uses four links, each with
Casimir 3/4, so \(H_0W_P=3W_P\) and \(H_0W_Q=3W_Q\). Differentiating the
eigen-equation at xi=0, including its scalar eigenvalue, now gives \[
 \phi_\xi=1+\frac\xi3(W_P+W_Q)+O(\xi^2),\qquad
 \log\rho_\xi=\frac{2\xi}3(W_P+W_Q)+O(\xi^2).
\] The remainders hold in every fixed Ck norm. To justify this
regularity, start with the L2 analytic Riesz vector. The eigen-equation
expresses H0 applied to that vector as multiplication by a fixed smooth
function times the vector, plus its scalar eigenvalue times the vector.
The elliptic estimate for the product Laplacian gains two Sobolev
derivatives at each iteration. Multiplication by the fixed Wilson
polynomial is bounded in each such Sobolev norm. Iterating also on the
differentiated equations gives analytic dependence in each fixed Sobolev
norm; Sobolev embedding then gives the stated fixed Ck bounds. Strict
positivity near 1 permits the ordinary logarithm with these same local
analytic bounds.

The normalization factor is \(Z_\xi=1+2\xi^2/9+O(\xi^3)\): the linear
coefficient has zero mean, its squared norm is 2/9, and every higher
coefficient of phi has zero mean by the chosen representative. This
records rather than hides the first state-normalization correction.

\subsection{3. Eliminated link and the true conditional
score}\label{eliminated-link-and-the-true-conditional-score}

Let S consist of one outer edge of P which is not in Q, and let R
contain the other six edges. Use lower-case s and r for their
configuration variables. Define the actual marginal and lift \[
 \rho_{R,\xi}(r)=\int\rho_\xi(s,r)d\lambda_S(s),\qquad Jf(s,r)=f(r).
\] J is an isometry from \(L^2(\rho_{R,\xi}d\lambda_R)\) into
\(L^2(\mu_\xi)\). Its adjoint is the literal conditional integral \[
 (J_\xi^*u)(r)=\rho_{R,\xi}(r)^{-1}
             \int u(s,r)\rho_\xi(s,r)d\lambda_S(s).
\] Put \(P_\xi=JJ_\xi^*\) and \(Q_\xi=I-P_\xi\). These symbols denote
orthogonal Hilbert-space projections, not the two plaquette holonomies.
In particular the formula specifies every dependence of the projection
on the true vacuum.

The ground-state transformed nonnegative operator is \[
 L_\xi=\psi_\xi^{-1}(H-E_\xi)\psi_\xi
       =-\kappa\sum_{f,a}\rho_\xi^{-1}X_{f,a}(\rho_\xi X_{f,a}).
\] For a smooth retained function f, direct differentiation and
conditional integration give \[
 L_\xi Jf=JA_\xi f+B_\xi f,
\] \[
 A_\xi f=-\kappa\sum_{f'\in R,a}
       \rho_{R,\xi}^{-1}X_{f',a}(\rho_{R,\xi}X_{f',a}f),
\] \[
 B_\xi f=-\kappa\sum_{f'\in R,a}
 (X_{f',a}\log\rho_\xi-X_{f',a}\log\rho_{R,\xi})X_{f',a}f.
\] The last function has zero conditional mean, since differentiating
the marginal integral gives
\(J_\xi^*(X\log\rho_\xi)=X\log\rho_{R,\xi}\). Thus it is the actual
component in the eliminated subspace. No independent-link vacuum is
used.

Choose the gauge-invariant retained observable
\(f_\xi=W_Q-\mu_\xi(W_Q)\). Haar integration over S annihilates WP and
all its derivatives on R, while leaving WQ unchanged. Therefore \[
 B_\xi f_\xi=-\frac{2b}{3}\,G+O(b^2/\kappa),\qquad
 G:=\sum_{a=1}^3(X_{e,a}W_P)(X_{e,a}W_Q).
\] Only the common edge contributes: all other derivatives of the two
plaquettes have disjoint supports. The remainder is in any fixed Sobolev
norm on this finite compact configuration space, at fixed kappa.

\subsection{4. Exact color contraction and two eliminated
eigencomponents}\label{exact-color-contraction-and-two-eliminated-eigencomponents}

Set \(Q'=BU_e^{-1}\). Then trace Q'=trace Q and \[
 X_{e,a}W_P=\operatorname{tr}(T_aP),\qquad
 X_{e,a}W_Q=-\operatorname{tr}(T_aQ').
\] The SU(2) Fierz identity in this exact basis is \[
 \sum_a\operatorname{tr}(T_aC)\operatorname{tr}(T_aD)
 =-\tfrac12[\operatorname{tr}(CD)-\tfrac12\operatorname{tr}C\operatorname{tr}D].
\] It follows by substituting the three Pauli matrices, or equivalently
their entry identity \(\sum_a(T_a)_{ij}(T_a)_{kl}
=-\tfrac12(\delta_{il}\delta_{jk}-\tfrac12\delta_{ij}\delta_{kl})\). Let
\[
 R_0=\operatorname{tr}(AB),\quad D_0=W_PW_Q,\quad
 D_1=D_0-\tfrac12R_0.
\] R0 is the outer six-edge Wilson loop and is independent of the shared
edge. The contraction, with its opposite shared-edge orientation
retained, is \[
 G=\tfrac12R_0-\tfrac14D_0
   =\tfrac38R_0-\tfrac14D_1.
\] To check all Haar inner products explicitly, for fixed Ue the
products P and Q' are independent Haar SU(2) matrices because A and B
are products of disjoint independent Haar edges. Write
\(P=p_0I+i\sum p_a\sigma_a\), \(Q'=q_0I+i\sum q_a\sigma_a\). Each p or q
is uniform on the unit three-sphere, so
\(\mathbb E p_\alpha p_\beta=\delta_{\alpha\beta}/4\), and the same
holds for q. Direct substitution gives \[
 G=-\sum_{a=1}^3p_aq_a,\quad
 R_0=2(p_0q_0-\sum_ap_aq_a),\quad
 D_1=3p_0q_0+\sum_ap_aq_a.
\] Consequently \[
 \|R_0\|_\lambda^2=1,\qquad \|D_1\|_\lambda^2=\tfrac34,
 \quad\langle R_0,D_1\rangle_\lambda=0,
 \quad\|G\|_\lambda^2=\tfrac3{16}.
\] Each of the six outer-link Laplacians acts on R0 and D1 by -3/4. For
the common edge, the product rule gives \[
 \Delta_e^T D_0=-\tfrac32D_0+2G=R_0-2D_0=-2D_1,
 \qquad \Delta_e^T R_0=0.
\] Therefore \[
 \kappa H_0R_0=\tfrac{9\kappa}{2}R_0,\qquad
 \kappa H_0D_1=\tfrac{13\kappa}{2}D_1.
\] Both have zero conditional mean over S, so these are also their
eigenvalues in the eliminated subspace at xi=0. In particular \[
 B_\xi f_\xi=-\frac b4R_0+\frac b6D_1+O(b^2/\kappa),\qquad
 \|B_\xi f_\xi\|_{\mu_\xi}^2
      =\frac{b^2}{12}+O(b^3/\kappa).
\] This proves that the eliminated quantum coupling on an actual
retained Wilson observable is nonzero for all sufficiently small
positive b/kappa. It is not a nonzero coupling assumed for an
unspecified observable.

\subsection{5. A full-system trial state with a strictly lower
excitation
energy}\label{a-full-system-trial-state-with-a-strictly-lower-excitation-energy}

Let Cxi be the nonnegative self-adjoint operator associated to the
closed form of Lxi restricted to \(\ker J_\xi^*\). The form domain is
the intersection of that closed subspace with H1. The smooth positive
density and its derivatives are bounded above and below on the compact
product. The explicit conditional integral above is a bounded map on H1;
therefore the domain is closed in the form norm and dense in the
eliminated subspace, by applying Qxi to smooth approximations. This
constructs Cxi without assuming a domain for a formal operator product.

For z\textgreater0, define \[
 w_{\xi,z}=(C_\xi+z)^{-1}B_\xi f_\xi,\qquad
 u_{\xi,z}=\psi_\xi(Jf_\xi-w_{\xi,z}).
\] All conditional projections, operators and vectors here commute with
the vertex gauge action; hence u is a physical gauge-invariant state.
Its inner product with the vacuum is zero, since f is centered and w has
zero conditional mean. Put \[
 d_\xi=\|f_\xi\|_{\rho_{R,\xi}}^2,\quad
 a_\xi=\kappa\sum_{f'\in R,a}\mu_\xi(|X_{f',a}f_\xi|^2),
 \quad I_{\xi,z}=\langle B_\xi f_\xi,w_{\xi,z}\rangle,
 \quad J_{\xi,z}=\|w_{\xi,z}\|^2.
\] I is real and nonnegative. Testing the resolvent equation by w gives
\(\mathfrak c_\xi(w,w)=I_{\xi,z}-zJ_{\xi,z}\). The cross form between Jf
and w is exactly \(\langle B_\xi f,w\rangle\). Expanding the full form
and norm, with their cross terms, gives the exact two Rayleigh values \[
 \mathcal R_{\rm retained}=\frac{a_\xi}{d_\xi},\qquad
 \mathcal R_{\rm full}(z)=
 \frac{a_\xi-I_{\xi,z}-zJ_{\xi,z}}{d_\xi+J_{\xi,z}},
\] and the exact decrease \[
 \mathcal R_{\rm retained}-\mathcal R_{\rm full}(z)
 =\frac{I_{\xi,z}+(z+\mathcal R_{\rm retained})J_{\xi,z}}
        {d_\xi+J_{\xi,z}}>0
\] whenever Bxi fxi is nonzero. This is a constructed improved vector in
the original seven-link interacting Hilbert space, not an assignment of
an effective classical potential to a quantum energy.

For completeness the expansions of this resolvent as xi tends to zero
are controlled on the fixed compact complex. The densities and their
first derivatives differ from their xi=0 values by O(xi), uniformly. The
conditional integral formula implies \(P_\xi-P_0=O(\xi)\) both on L2 and
on H1. All Pxi have exactly the same range, the functions of the
retained links. Thus \(P_\xi P_0=P_0\), \(P_0P_\xi=P_\xi\), and
consequently \(Q_0Q_\xi=Q_0\), \(Q_\xi Q_0=Q_\xi\). On \(\ker P_0\), the
map \(T_\xi=Q_\xi|_{\ker P_0}\) therefore has the exact inverse
\(Q_0|_{\ker P_\xi}\). Both preserve H1 by the conditional integral
bounds. This is an exact identification of the moving form domains, not
an assumption that the conditional projections are independent of xi.
After this identification, the forms Cxi+z are uniformly coercive in H1
and differ from the xi=0 form by O(xi) in H1 form norm. Subtracting
their weak resolvent equations and using coercivity proves an O(xi)
difference of solution maps from the dual of H1 to H1. Applying this to
the already proved \(B_\xi f_\xi=b(-R_0/4+D_1/6)+O(b\xi)\) justifies,
without an unproved interchange of spectra or projections, \[
 w_{\xi,z}=-\frac{b}{4(9\kappa/2+z)}R_0
             +\frac{b}{6(13\kappa/2+z)}D_1+O(\xi^2)
\] at fixed positive kappa and z. The error is in H1 and includes the
variation of the conditional projection and state density.

In particular, with the retained choice z=kappa, \[
 I_{\xi,\kappa}=\frac{b^2}{\kappa}
           \left(\frac1{88}+\frac1{360}\right)+O(b^3/\kappa^2)
           =\frac{7b^2}{495\kappa}+O(b^3/\kappa^2),
\] \[
 J_{\xi,\kappa}=\frac{b^2}{\kappa^2}
           \left(\frac1{484}+\frac1{2700}\right)+O(b^3/\kappa^3)
           =\frac{199b^2}{81675\kappa^2}+O(b^3/\kappa^3).
\] Haar integration gives \(d_0=1\) and \(a_0=3\kappa\); their changes
are O(xi) with the respective retained scales. Substitution in the exact
decrease formula therefore proves the explicit positive coefficient \[
 \boxed{\mathcal R_{\rm retained}-\mathcal R_{\rm full}(\kappa)
     =\frac{1951}{81675}\frac{b^2}{\kappa}
          +O(b^3/\kappa^2).}
\] In the original coupling and spatial spacing this coefficient is
\(1951/[653400\,g_{\rm YM}^6a]\), with remainder of order
\(1/(g_{\rm YM}^{10}a)\) in the stated strong-coupling regime.

\subsubsection{5.1 Absolute excitation energies, not only their
difference}\label{absolute-excitation-energies-not-only-their-difference}

The absolute trial-state energies can also be computed to this order.
Let \(F=W_P+W_Q\), and write the mean-one vacuum representative as
\(\phi_\xi=1+\xi F/3+\xi^2\phi_2+O(\xi^3)\). Its eigen-equation and mean
constraint give \[
 H_0\phi_2=(F^2-2)/3,\qquad \int\phi_2d\lambda=0.
\] The subtracted 2 is the exact \(\int F^2d\lambda\), not a removed
plaquette. The identity \(W_Q^2=1+\chi_{\rm ad}(Q)\) follows by
\(2\cos\theta\) squared, with \(\chi_{\rm ad}=1+2\cos2\theta\). The
adjoint Casimir is 2 on each of Q's four links, so
\(H_0\chi_{\rm ad}(Q)=8\chi_{\rm ad}(Q)\). The independent Haar
holonomies P and Q satisfy \[
 \int W_Q^2=1,\qquad\int W_Q^4=2,\qquad
 \int F^2W_Q^2=3,\qquad\int F^2\chi_{\rm ad}(Q)=1.
\] For the fourth moment, \(W_Q=2q_0\) and the S3 marginal used above
gives \(\int q_0^4=1/8\), either by its beta integral or the
corresponding even moment recurrence. Cross terms vanish by Haar
symmetry of P and Q. Self-adjointness of H0 now yields \[
 \int\phi_2 W_Q^2d\lambda
 =\int\phi_2\chi_{\rm ad}(Q)d\lambda=\frac1{24}.
\] Using the full state-density expansion, including Zxi, \[
 \rho_\xi=1+\frac{2\xi}3F
     +\xi^2\left(\frac{F^2}{9}+2\phi_2-\frac29\right)+O(\xi^3),
\] therefore gives \[
 r_\xi:=\mu_\xi(W_Q^2)=1+\frac7{36}\xi^2+O(\xi^3),\qquad
 \mu_\xi(W_Q)=\frac23\xi+O(\xi^2),
\] \[
 d_\xi=r_\xi-\mu_\xi(W_Q)^2=1-\frac14\xi^2+O(\xi^3).
\] For a single edge of the SU(2) loop,
\(\sum_a|X_{e,a}W_Q|^2=(4-W_Q^2)/4\), by the displayed Pauli
coordinates. Summing its four edges gives exactly
\(a_\xi=\kappa(4-r_\xi)\). Hence the absolute excitation Rayleigh values
relative to the same actual vacuum are \[
 \boxed{\mathcal R_{\rm retained}
    =3\kappa+\frac59\frac{b^2}{\kappa}+O(b^3/\kappa^2),}
\] \[
 \boxed{\mathcal R_{\rm full}(\kappa)
    =3\kappa+\frac{43424}{81675}\frac{b^2}{\kappa}
         +O(b^3/\kappa^2).}
\] The latter coefficient is precisely \(5/9-1951/81675=43424/81675\).
The elimination correction decreases the variational energy at the same
positive coupling, but both trial energies increase from their common
zero-b value to this order. Neither the decrease nor the positive
second-order coefficients determine the full spectrum or a limiting
four-dimensional gap.

For completeness the energy origin is explicit as well. Taking the Haar
mean of the second-order eigen-equation gives
\(e(\xi)=-2\xi^2/3+O(\xi^3)\). The actual, unsubtracted Hamiltonian
therefore has vacuum energy \[
 E_\xi=4b-\frac{2b^2}{3\kappa}+O(b^3/\kappa^2).
\] Adding this same scalar to each excitation quotient gives the two
trial quotients for H itself: \[
 3\kappa+4b-\frac19\frac{b^2}{\kappa}+O(b^3/\kappa^2),\qquad
 3\kappa+4b-\frac{11026}{81675}\frac{b^2}{\kappa}
       +O(b^3/\kappa^2),
\] respectively. The mass-gap question concerns energies above the
vacuum, so the preceding excitation quotients, not this common scalar
shift, are the quantities used in the variational comparison.

\subsection{6. Meaning and sources}\label{meaning-and-sources}

The correction is already present in a nonlinear gauge-invariant quantum
calculation: it comes from covariance between a retained loop and the
eliminated link through their common plaquette edge. Its two explicit
Casimir channels, 9kappa/2 and 13kappa/2, supply the coefficient above.
The improved state still has an excitation Rayleigh value near 3kappa in
this regime. No small-energy continuum state or four-dimensional quantum
mass-gap disproof is inferred from this finite-complex result.

The exact general marginal, resolvent and trial-state identities are
proved in the companion source
\texttt{spatial\_vacuum\_elimination.tex}; the calculation above
evaluates its nonzero coupling rather than assuming it. The original
coefficient dictionary is proved from the temporal Wilson integral in
\texttt{temporal\_transfer.tex}. The compact-ground-state and Fierz
identities are expanded in \texttt{wilson\_finite\_lattice.tex}.

The general elimination method is established literature, not claimed as
a new discovery: G. Dusson, I. M. Sigal and B. Stamm, \emph{The
Feshbach--Schur map and perturbation theory}, pp.~65--88 (2021),
\url{https://doi.org/10.4171/ECR/18-1/5},
\url{https://arxiv.org/abs/2105.02058}. The lattice Hamiltonian context
is J. Kogut and L. Susskind, \emph{Hamiltonian formulation of Wilson's
lattice gauge theories}, Physical Review D 11 (1975), 395--408,
\url{https://doi.org/10.1103/PhysRevD.11.395}.

\end{document}
